\documentclass[12pt]{article}
\usepackage[margin=1in]{geometry}
\usepackage{booktabs,makecell,threeparttable}

\begin{document}

\providecommand{\StructuralRobustnessPath}{appendix/structural-robustness}

\subsection{The Geographic Support of the Policy Exercise}
\label{app:policy-distance-cap}

Our maintained policy exercise permits assignments up to 3.5 km. A seemingly
natural alternative would cap assignments at 2.5 km, near the upper end of the
experimental distance distribution. In this setting, however, a 2.5 km cap is
not an informative structural-policy robustness check because it nearly
recreates the geographic separation built into the experimental design.

The experimental clusters were constructed from 2.5 km inner buffers around
candidate points of treatment. The selection algorithm used 3 km outer radii
for the Control/Ink group and 4 km outer radii for the Bracelet/Calendar group,
and removed overlapping pieces of candidate cluster areas. This design helped
separate experimental clusters, but it also means that distinct experimental
communities rarely share a feasible distribution site within 2.5 km.

This limitation does not arise because the policy exercise retained only
schools from a 2.5 or 3 km search radius. Its candidate-site file contains all
1,451 documented schools and was generated using an unrestricted
candidate-distance cutoff. The binding restriction is instead the overlap of the villages'
feasible-site neighborhoods. Table~\ref{tab:policy-distance-cap-feasibility}
shows this directly. At 2.5 km, only 53 candidate sites are feasible for more
than one village. Even if demand were identical at every site and the planner
only minimized the number of sites, the exact geographic minimum would be 114
sites for 144 villages. Both policy regimes attain this same floor. The zero
Control--Bracelet difference at 2.5 km therefore follows mechanically from the
experimental geography rather than from a failure of the social-image demand
mechanism.

\begin{table}[htbp]
  \centering
  \small
  \begin{threeparttable}
    \caption{Distance caps and the geographic feasibility of site consolidation}
    \label{tab:policy-distance-cap-feasibility}
    \begin{tabular}{rrrrrr}
\toprule
Max. distance & Feasible links & Shareable sites & \makecell{Geographic\\minimum} & \makecell{Control\\sites} & \makecell{Bracelet\\sites} \\
\midrule
2.50 km & 914 & 53 & 114 & 114 & 114 \\
2.75 km & 1,085 & 109 & 99 & 108 & 99 \\
3.00 km & 1,256 & 186 & 89 & 106 & 89 \\
3.25 km & 1,430 & 272 & 81 & 105 & 81 \\
3.50 km & 1,639 & 398 & 71 & 101 & 76 \\
\bottomrule
\end{tabular}

    \begin{tablenotes}[flushleft]
      \footnotesize
      \item \textit{Notes:} ``Feasible links'' counts village--candidate-site
      pairs within the cap; ``shareable sites'' are feasible for at least two
      villages. ``Geographic minimum'' is the exact site-count minimum with
      identical demand. Control and Bracelet preserve population-weighted
      experimental-allocation take-up using the component-wise posterior
      median of 1,600 retained baseline draws. The exercise covers 144
      villages and 1,451 candidate schools.
    \end{tablenotes}
  \end{threeparttable}
\end{table}

Feasible consolidation rises rapidly just beyond the design radius. At 3
km---close to the maximum experimental assignment distance of 2.85
km---Bracelet observability saves 17 sites. At the maintained 3.5 km cap, the
geographic minimum is 71 sites and the structural allocations use 101 sites
under Control observability versus 76 under Bracelet observability, saving 25.
We retain 3.5 km as a modest spatial extrapolation needed for cross-community
consolidation. The 2.5 km result is a design diagnostic, not evidence against
the policy mechanism.

\end{document}
